% !TeX root = ../../Book2.tex
% 纲要标题：### A.2 主理想整环的计算与一般环的障碍

% 纲要要点（供目录核对）：
% * 主理想整环上的计算和一般环上的障碍
% #### A.2.1 有效表示与终止条件
% #### A.2.2 非主理想造成的障碍
% #### A.2.3 系数增长与模素数检验

\section{主理想整环的计算与一般环的障碍}
\label{b2:sec:A02}

Smith 理论把主理想条件转化为结构结论，但计算还取决于环如何给出。这里先区分“存在所需元素”与“能够求出所需元素”，再用一个二元多项式矩阵说明：离开主理想整环后，失败的可能是标准形本身，而不只是现有算法太慢。

\subsection{有效表示与终止条件}
\label{b2:subsec:A02:01}

对给定主理想整环 \(R\)，执行第四章的 Bézout 约化至少需要以下操作：有限地表示输入元素，精确执行环运算和相等判定；判定 \(a\mid b\) 并在成立时求出商；对非同时为零的 \(a,b\)，求出
\[
g=ua+vb,\qquad a=ga',\qquad b=gb',
\]
其中 \(g\) 生成 \((a,b)\)；还需判定单位并求其逆。若要求唯一的元素输出，须进一步指定可执行的相伴类归一化。这里所说的有效计算条件是对环的表示及操作提出的要求，并不是主理想整环定义中额外隐含的一部分。

\ProseParagraphBreak
在这些操作已给出的前提下，矩阵
\[
E=\begin{pmatrix}u&v\\-b'&a'\end{pmatrix}
\]
满足 \(\det E=ua'+vb'=1\)，并把 \((a,b)^{\mathsf T}\) 化为 \((g,0)^{\mathsf T}\)。若 \(a\nmid b\)，主元理想从 \((a)\) 严格扩大为 \((g)\)。\cref{b2:lem:pid-ascending-chain}已经证明主理想整环中的理想升链最终稳定，因此反复使用这样的约化也会终止，不必预先找到 Euclid 函数。这个论证保证给定输入在有限步后结束，却没有给出随输入长度变化的统一运行时间界。

\ProseParagraphBreak
这里说域 \(k\) 有效给定，是指它的元素有有限表示，且加、减、乘、非零元的除法以及相等和零判定都能在有限步内执行。整数和这样的域上的一元多项式环满足上述计算条件。例如，多项式扩展 Euclid 算法不仅输出最大公因式，还逐步保存 Bézout 系数；首一归一化只需除以一个非零域元素。这一过程及次数估计可参阅 Shoup 的《A Computational Introduction to Number Theory and Algebra》\cite[Section 17.3, Theorems 17.2, 17.4, pp. 469--472]{Shoup2008NumberTheoryAlgebra}。域的有效性不能省略：给出一个任意实数的近似值，并不同时给出了精确的零判定。

\subsection{非主理想造成的障碍}
\label{b2:subsec:A02:02}

\begin{proposition}\label{b2:prop:A-nonprincipal-obstruction}
设 \(k\) 为域，\(R=k[x,y]\)。行矩阵 \(A=(x\ y)\) 不与任何 \((d\ 0)\) 在 \(R\) 上幺模等价。
\end{proposition}
\begin{proof}
可逆行列变换保持所有条目生成的理想。事实上，变换后每个条目是原条目的 \(R\) 线性组合，给出一个包含关系；用逆变换得到反向包含。因此若这样的等价存在，便有 \((x,y)=(d)\)。

由此 \(d\) 同时整除 \(x,y\)。写 \(x=dh\)，并把等式看作 \(k[x][y]\) 中关于 \(y\) 的多项式等式。由于 \(x\ne0\) 且非零多项式的 \(y\) 次数在乘法下相加，有 \(\deg_y d=0\)，即 \(d\in k[x]\)。再写 \(y=dg\)，比较 \(y\) 的一次项系数，得到 \(d\) 在 \(k[x]\) 中整除 \(1\)，故 \(d\) 是单位。这会给出 \((x,y)=R\)。但代入 \(x=y=0\)，每个 \(xf+yg\) 都取零值，不可能等于一，矛盾。
\end{proof}

这个例子在域 \(k(x,y)\) 上只有秩一，仍不能在原环上化成一个非零对角元。秩丢失了非主理想 \((x,y)\) 的信息。一般交换环上的条目理想和高阶子式理想仍有意义，但不能预设每个这样的理想都有单个生成元；第三卷第12.3节有关模展示与 Fitting 理想的讨论将继续利用这些不变量。

\ProseParagraphBreak
有零因子时还要另查核的计算。例如在 \(R=\mathbb Z/6\mathbb Z\) 中，对角方程 \(2z=0\) 有两个解 \(z=0,3\)。因而即使矩阵已经对角化，也不能照用\cref{b2:prop:A-smith-solving}中“非零对角元的对应齐次坐标为零”这一步；该步明确使用了整环的非零消去律。

\subsection{系数增长与模素数检验}
\label{b2:subsec:A02:03}

整数运算的代价不仅取决于运算次数，还取决于中间整数的位数。行列相减虽不会引入分母，倍数的累积仍可能使换基矩阵的条目很大。对有理数先清除分母，可以把零空间或线性方程化为整数计算；若随后研究整数余核，却必须保留这一缩放的含义，因为乘一个非单位会改变整数关系模。

\begin{proposition}\label{b2:prop:A-modular-rank}
设 \(A\in M_{m\times n}(\mathbb Z)\)，\(p\) 是素数，\(\overline A\in M_{m\times n}(\mathbb F_p)\) 是将 \(A\) 的条目逐一模 \(p\) 所得的矩阵，则
\[
\operatorname{rank}_{\mathbb F_p}\overline A\le
\operatorname{rank}_{\mathbb Q}A.
\]
若某个 \(r\) 阶子式在模 \(p\) 后非零，则原矩阵的有理数秩至少为 \(r\)。若同时给出了 \(A=UV\)，其中 \(U\) 为 \(m\times r\) 整数矩阵、\(V\) 为 \(r\times n\) 整数矩阵，那么原矩阵的有理数秩恰为 \(r\)。
\end{proposition}
\begin{proof}
\(\overline A\) 的非零 \(r\) 阶子式来自一个不被 \(p\) 整除的整数子式，所以该整数非零，原矩阵秩至少为 \(r\)。取 \(r\) 为模 \(p\) 的秩即得不等式。另一方面，若 \(A=UV\)，它在 \(\mathbb Q\) 上的像包含于 \(U\) 的至多 \(r\) 维的列空间，所以秩至多为 \(r\)，与前面的下界合并得到等号。
\end{proof}

例如 \(\operatorname{diag}(1,6)\) 在 \(\mathbb Q\) 上秩为二，模二或模三后秩降为一，模五后秩仍为二。模 \(p\) 为零的整数子式未必在整数中为零；模 \(p\) 非零的子式则能证明对应的整数子式非零。

\begin{proposition}\label{b2:prop:A-bounded-congruence}
设 \(B\ge0\) 为实数，\(s\ge1\) 为整数，\(z\in\mathbb Z\)、\(|z|\le B\)，并设两两互素的正整数 \(q_1,\ldots,q_s\) 满足
\[
z\equiv0\pmod {q_i}\quad(1\le i\le s),\qquad
q_1\cdots q_s>B.
\]
则 \(z=0\)。若 \(z,z'\in[-B,B]\cap\mathbb Z\) 在这些模数下逐一同余，且 \(q_1\cdots q_s>2B\)，则 \(z=z'\)。
\end{proposition}
\begin{proof}
两两互素性说明乘积 \(M=q_1\cdots q_s\) 整除 \(z\)。非零的 \(M\) 倍数绝对值至少为 \(M\)，与 \(|z|\le B<M\) 矛盾。第二个结论对 \(z-z'\) 使用第一个结论，因为 \(|z-z'|\le2B<M\)。
\end{proof}

这使模运算也能提供确定性的等式证书。检验整数矩阵的等式 \(PAQ=D\) 时，可以直接从输入条目取得上界，无须先算出整个乘积。设 \(m,n\geq1\)，\(A,D\in M_{m\times n}(\mathbb Z)\)、\(P\in M_m(\mathbb Z)\)、\(Q\in M_n(\mathbb Z)\)，并记 \(\lVert M\rVert_{\max}=\max_{i,j}\lvert M_{ij}\rvert\)。每个乘积条目是 \(mn\) 项的和，因而
\[
\begin{aligned}
B&\coloneqq mn\lVert P\rVert_{\max}\cdot\lVert A\rVert_{\max}
\cdot\lVert Q\rVert_{\max}+\lVert D\rVert_{\max},\\
\lVert PAQ-D\rVert_{\max}&\leq B.
\end{aligned}
\]
这个粗界只需计算四个输入矩阵的最大条目绝对值。选取乘积大于 \(B\) 的两两互素模数，并在每个模数下验证 \(PAQ=D\)，\cref{b2:prop:A-bounded-congruence}就逐条目保证整数等式成立。

\ProseParagraphBreak
没有上界时，即使多次模检验都给出零，也不能断定整数等式成立：模数的乘积本身非零，却在每个模数下都为零。用中国剩余定理恢复有界整数的矩阵计算方法，见 Shoup 的上述著作\cite[Section 4.4, pp. 84--86]{Shoup2008NumberTheoryAlgebra}；这里仅使用有界整数的唯一性，不对本附录的 Smith 算法声称复杂度最优。
